% Materiał z AF 25; fragment do rozdziału o iloczynie induktywnym.
\subsection{Identyfikacje przestrzeni sprzężonych do iloczynów tensorowych}

\begin{stwierdzenie}[Zestawienie kanonicznych odwzorowań]
\label{prop:summary25}
Dla każdej liczby całkowitej $d\ge1$ i ustalonego
ciała $\K\in\{\R,\C\}$ zachodzi schemat
\[
 \begin{array}{ccccc}
 \mathcal D(\R^d)&\longrightarrow&\mathcal S(\R^d)
 &\longrightarrow&\mathcal E(\R^d)\\
 \mathcal D'(\R^d)&\longleftarrow&\mathcal S'(\R^d)
 &\longleftarrow&\mathcal E'(\R^d).
 \end{array}
\]
W pierwszym wierszu strzałki oznaczają ciągłe
różnowartościowe odwzorowania liniowe o gęstych obrazach,
w drugim zaś ich transpozycje, ciągłe i różnowartościowe
zarówno w topologiach słabych-*, jak i silnych.
Żadna ze strzałek nie jest izomorfizmem topologicznym
na obraz z topologią podprzestrzeni.

Dla dodatnich całkowitych $d_1,d_2$ oraz niepustych
zbiorów otwartych $\Omega\subseteq\R^{d_1}$,
$\Lambda\subseteq\R^{d_2}$ istnieją kanoniczne
izomorfizmy liniowo-topologiczne
\[
\begin{aligned}
 \mathcal E(\Omega)\widehat\otimes_\pi\mathcal E(\Lambda)
 &\simeq\mathcal E(\Omega\times\Lambda),\\
 \mathcal S(\R^{d_1})\widehat\otimes_\pi\mathcal S(\R^{d_2})
 &\simeq\mathcal S(\R^{d_1+d_2}),\\
 \mathcal D(\Omega)\widehat\otimes_i\mathcal D(\Lambda)
 &\simeq\mathcal D(\Omega\times\Lambda).
\end{aligned}
\]
Każdy izomorfizm przedłuża odwzorowanie
$f\otimes g\mapsto ((x,y)\mapsto f(x)g(y))$.
W pierwszych dwóch wzorach można zastąpić $\pi$
przez $\varepsilon$. W trzecim wzorze nie wolno
ogólnie zastąpić $i$ przez $\pi$, choć na algebraicznym
iloczynie tensorowym przestrzeni $\mathcal D$
zachodzi $\varepsilon=\pi$.
\end{stwierdzenie}
\begin{proof}
Pierwszy wiersz wynika z Twierdzenia~
\ref{thm:zanurzenia25}, drugi z Wniosku~
\ref{cor:dual-lancuch25}; różnicę topologii
na obrazach dualnych wykazano w Przykładzie~
\ref{ex:escaping-deltas25}. Pierwsze dwa izomorfizmy
oraz dopuszczalność $\varepsilon$ wynikają z
Twierdzenia~\ref{thm:smooth-schwartz-tensor24}.
Trzeci wynika z Twierdzenia~\ref{thm:D-i-product24};
równość $\varepsilon=\pi$ wynika z nuklearności
$\mathcal D$ (Twierdzenia~\ref{thm:nuclear-D23}
i~\ref{thm:nuclear-pi-eps24}). Nie można zastąpić
$i$ przez $\pi$ już dla $\Omega=\Lambda=\R$,
na mocy Uwagi~\ref{rem:D-not-projective24}.
\end{proof}

\begin{uwaga}[Odczytywanie schematu po stronie dualnej]
\label{rem:kernel-summary25}
Jeżeli $J:X\to Y$ jest izomorfizmem liniowo-topologicznym,
to jego transpozycja $J':Y'\to X'$ jest także
izomorfizmem dla obu topologii dualnych.
Istotnie, $J'$ jest ciągłe przez równość seminorm
silnych i słabych-* z dowodu Lematu~
\ref{lem:transpozycja25}, a jego odwrotnością jest
$(J^{-1})'$, ciągłe z tego samego powodu.
Dlatego transpozycje trzech izomorfizmów tensorowych
identyfikują przestrzenie dualne produktów
z odpowiednimi przestrzeniami dystrybucji na
iloczynie, z zachowaniem topologii słabej-* i silnej.

Konkretnie, w każdym przypadku funkcjonał $K$
na przestrzeni funkcji na iloczynie odpowiada
formie dwuliniowej
$b_K(f,g)=\langle K,f\otimes g\rangle$.
Dla $\mathcal E$ i $\mathcal S$ jest to bijekcja
liniowa między dystrybucjami z odpowiedniej
przestrzeni dualnej na iloczynie a \emph{łącznie}
ciągłymi formami dwuliniowymi na parze czynników.
Wynika to z własności uniwersalnej $\pi$
(Twierdzenie~\ref{thm:pi-universal24}) i jednoznacznego
przedłużenia funkcjonału na uzupełnienie.
Dla $\mathcal D$ analogiczna bijekcja dotyczy
\emph{osobno} ciągłych form, zgodnie z
Wnioskiem~\ref{cor:kernel-i24}.
W obu przypadkach jednoznaczność jądra wynika
z gęstości liniowej powłoki tensorów elementarnych
w odpowiednim uzupełnieniu.

Nie przypisujemy tutaj przestrzeni form dwuliniowych
dodatkowej topologii. Stwierdzenie o izomorfizmach
topologicznych dotyczy dualnych przestrzeni
uzupełnionych iloczynów tensorowych, nie dowolnie
wybranych topologii na przestrzeniach form.
\end{uwaga}

\begin{uwaga}[Zakres dalszych zastosowań]
\label{rem:future25}
Identyfikacje przestrzeni sprzężonych i twierdzenie o jądrze
opisują reprezentacje ciągłych operatorów na przestrzeniach
funkcji testowych. Zastosowania do konkretnych równań
różniczkowych rozwijamy w skrypcie RRC, a ogólną teorię
operatorów w analizie funkcjonalnej. W dalszej części
niniejszego skryptu przechodzimy do analizy harmonicznej.
\end{uwaga}
